\documentclass[11pt]{article}

% =========================================================
% Packages
% =========================================================
\usepackage[margin=1in]{geometry}
\usepackage{amsmath,amssymb,amsthm,mathtools}
\usepackage{booktabs}
\usepackage{tabularx}
\usepackage{array}
\usepackage{multirow}
\usepackage{enumitem}
\usepackage{microtype}
\usepackage{algorithm}
\usepackage{algpseudocode}
\usepackage{hyperref}
\usepackage{xcolor}

\hypersetup{
    colorlinks=true,
    citecolor=blue,
    linkcolor=blue,
    urlcolor=blue
}

\newtheorem{proposition}{Proposition}
\newtheorem{theorem}{Theorem}
\newtheorem{corollary}{Corollary}
\newtheorem{remark}{Remark}
\newtheorem{definition}{Definition}

\newcommand{\KL}{D_{\mathrm{KL}}}
\newcommand{\E}{\mathbb{E}}
\newcommand{\cL}{\mathcal{L}}
\newcommand{\cD}{\mathcal{D}}
\newcommand{\cR}{\mathcal{R}}
\newcommand{\sg}{\operatorname{sg}}
\newcommand{\relu}{[\cdot]_+}
\newcommand{\simcos}{\operatorname{cos}}
\newcommand{\R}{\mathbb{R}}

% =========================================================
% Title
% =========================================================

\title{
\textbf{Beyond Representation Alignment:}\\
Semantic Distortion-Rate Learning for Matryoshka Representations
}

\author{
Anonymous Proposal
}

\date{}

\begin{document}

\maketitle

% =========================================================
% Abstract
% =========================================================

\begin{abstract}
Matryoshka Representation Learning (MRL) enables a single embedding to operate at multiple dimensionalities by requiring prefixes of a high-dimensional representation to remain useful under truncation.
Despite its practical success, existing approaches primarily optimize prefix-specific task losses or align truncated representations with richer hidden representations.
Such objectives do not explicitly characterize \emph{which semantic information} should survive at each representation budget.

We propose \textbf{Semantic Distortion-Rate Matryoshka Representation Learning (SDR-MRL)}, a framework that views each Matryoshka prefix as a rate-constrained semantic representation.
Instead of maximizing mutual information between hidden states, SDR-MRL defines the distortion of a prefix through the discrepancy between the semantic neighborhood distribution induced by a full-dimensional teacher and the distribution recoverable from the prefix.
We show that this neighborhood distortion admits a variational decomposition into the conditional semantic information lost by the prefix and a decoder approximation term.
At the optimal semantic decoder, the reduction in distortion obtained by appending a new coordinate block is exactly the conditional semantic information contributed by that block given the preceding prefix.

This perspective connects semantic geometry and information theory without requiring explicit high-dimensional mutual-information estimation.
We further introduce a multi-rate training objective, an optional monotonic refinement regularizer, a semantic distortion-rate evaluation protocol, and a set of experiments designed to measure the \emph{price of nestedness} relative to independently optimized embeddings.
The central hypothesis is that directly preserving rate-conditioned semantic neighborhood geometry will substantially improve extreme low-dimensional retrieval while reducing the performance gap between nested and independently trained representations.
\end{abstract}

% =========================================================
\section{Introduction}
% =========================================================

Dense vector representations are now a central interface between pretrained models and downstream systems.
They support semantic search, retrieval-augmented generation, clustering, classification, and large-scale nearest-neighbor indexing.
However, the computational and storage cost of an embedding system grows with representation dimensionality.
A representation that is appropriate for a high-resource retrieval server may be unnecessarily expensive for an edge device, a memory-constrained index, or an application requiring billions of stored vectors.

Matryoshka Representation Learning (MRL) addresses this problem by learning a single representation
\begin{equation}
    h(x) \in \R^D
\end{equation}
whose prefixes
\begin{equation}
    h_{1:d_1}(x),\;
    h_{1:d_2}(x),\ldots,
    h_{1:d_K}(x),
    \qquad
    d_1 < d_2 < \cdots < d_K=D,
\end{equation}
remain useful at multiple deployment budgets~\cite{kusupati2022matryoshka}.
This produces an attractive ``train once, deploy at many capacities'' paradigm with no additional inference architecture.

The success of MRL nevertheless raises a more fundamental question:
\begin{quote}
\emph{What semantic information should be allocated to each successive prefix?}
\end{quote}

Standard MRL supervises multiple dimensionalities with task losses, but the objective does not explicitly define what information should be placed in early versus late coordinates.
Recent approaches improve this organization in different ways.
ESE introduces PCA-driven compression to move salient variation toward early dimensions~\cite{li2025ese}.
Matryoshka-Adaptor learns mappings that compress embeddings from pretrained or black-box models~\cite{yoon2024matryoshka}.
SMEC addresses gradient interference and adaptive dimension selection for retrieval compression~\cite{zhang2025smec}.
Most recently, MIPIC introduces Self-Distilled Intra-Relational Alignment (SIA), including top-$k$ CKA alignment between truncated and full hidden states, together with Progressive Information Chaining (PIC) across model depth~\cite{huy2026mipic}.

These methods substantially strengthen Matryoshka training, but they largely approach the problem through \emph{representation supervision}: a compressed representation should agree with a richer representation, a principal subspace, or an adjacent hidden checkpoint.
We argue that representation agreement and semantic information preservation are related but distinct objectives.

Consider two full-dimensional embeddings related by an orthogonal transformation.
Their full-space Euclidean geometry and inner products can be identical, yet coordinate truncation can produce dramatically different low-dimensional representations.
A semantic direction concentrated in the first few coordinates may be spread across all coordinates by the transformation.
Consequently,
\begin{equation}
\boxed{
\text{equivalent full-space geometry}
\;\not\Rightarrow\;
\text{equivalent prefix utility}.
}
\end{equation}
The Matryoshka problem therefore contains an \emph{ordered information-allocation problem} that is absent from ordinary fixed-dimensional representation learning.

We propose to model this problem as \textbf{semantic successive refinement}.
Each prefix is treated as a representation available at a particular operational storage rate.
Rather than asking whether a compressed hidden state resembles the full hidden state, we ask:
\begin{quote}
\emph{How much of the semantic neighborhood structure remains recoverable at a given embedding rate?}
\end{quote}

To answer this question, we define a semantic random variable $S$ representing the neighborhood relation induced by a teacher representation.
For a Matryoshka prefix $Z_k$, we measure the discrepancy between the teacher semantic distribution $p_T(S|X)$ and a prefix-induced decoder $q_k(S|Z_k)$.
The resulting semantic neighborhood distortion is
\begin{equation}
    \cD_k
    =
    \E_X
    \KL
    \left(
        p_T(S|X)
        \,\Vert\,
        q_k(S|Z_k)
    \right).
\end{equation}

Our key observation is that this apparently geometric quantity has a direct information-theoretic interpretation:
\begin{equation}
\boxed{
\cD_k
=
I_T(S;X|Z_k)
+
\epsilon_k,
}
\end{equation}
where $\epsilon_k$ is a semantic-decoder approximation gap.
Thus, semantic neighborhood distortion is a variational upper bound on the semantic information that remains unavailable after compression.
With an optimal semantic decoder,
\begin{equation}
    \cD_k^\star = I_T(S;X|Z_k).
\end{equation}

Moreover, because Matryoshka representations are nested,
\begin{equation}
    Z_k=(Z_{k-1},B_k),
\end{equation}
where $B_k$ denotes the newly appended coordinate block.
At the optimal decoder,
\begin{equation}
\boxed{
\cD_{k-1}^\star-\cD_k^\star
=
I_T(S;B_k|Z_{k-1}).
}
\end{equation}
The improvement obtained by increasing dimensionality therefore has a precise interpretation: it is the conditional semantic contribution of the newly available block.

This leads to SDR-MRL, which optimizes semantic distortion across a distribution of deployment rates rather than explicitly maximizing mutual information between high-dimensional hidden representations.

\paragraph{Contributions.}
This proposal investigates four main contributions:

\begin{enumerate}[leftmargin=1.5em]
    \item \textbf{Semantic distortion formulation.}
    We define rate-conditioned semantic neighborhood distortion for nested embeddings and connect it to conditional semantic information through a variational decomposition.

    \item \textbf{Successive refinement interpretation.}
    We show that, at an optimal semantic decoder, the distortion reduction between adjacent prefixes equals the conditional semantic information contributed by the newly added dimensions.

    \item \textbf{Multi-rate Matryoshka objective.}
    We introduce an objective that minimizes expected semantic distortion under a deployment-rate distribution, together with an optional monotonic refinement regularizer.

    \item \textbf{Evaluation beyond per-dimension accuracy.}
    We propose semantic distortion-rate curves, neighborhood preservation, and the price of nestedness relative to independently trained models as diagnostics of information organization.
\end{enumerate}


% =========================================================
\section{Related Work}
% =========================================================

\subsection{Matryoshka Representation Learning}

MRL learns nested representations by applying supervision at multiple prefix dimensionalities~\cite{kusupati2022matryoshka}.
The resulting embeddings encode information at multiple granularities and allow the deployment dimensionality to be selected after training.
The original formulation is deliberately simple and broadly applicable, but its multi-prefix objective does not explicitly characterize semantic information allocation across coordinates.

Subsequent work has explored several complementary directions.
Matryoshka-Adaptor learns post-hoc transformations for compressing existing embeddings and is applicable even to black-box embedding APIs~\cite{yoon2024matryoshka}.
ESE jointly targets model-depth and embedding-dimension elasticity; its learn-to-compress component uses PCA-related supervision to compact salient features into early dimensions~\cite{li2025ese}.
SMEC studies retrieval-oriented Matryoshka compression, identifying gradient instability and information degradation as important practical issues and introducing sequential optimization, adaptive dimension selection, and cross-batch memory~\cite{zhang2025smec}.

These approaches improve practical compression but leave open a general criterion for determining whether a semantic structure is placed at an information-efficient point in the nested representation.

\subsection{MIPIC and Cross-Dimensional Alignment}

MIPIC explicitly targets information organization across both embedding dimensionality and network depth~\cite{huy2026mipic}.
Its Self-Distilled Intra-Relational Alignment component encourages truncated representations to preserve token-level structure from full-dimensional states, including attention matching and top-$k$ CKA-based alignment.
Its Progressive Information Chaining component transfers semantic information through a sequence of layer--dimension checkpoints.

MIPIC is particularly relevant to our work because it recognizes that independent prefix supervision alone does not sufficiently control information organization.
Our approach differs in the object being preserved.
MIPIC primarily aligns hidden-state structure and progressively connects intermediate representations.
SDR-MRL instead defines a semantic random variable over \emph{inter-example neighborhoods} and asks how much of this semantic relation is recoverable from each prefix.

The distinction is important for retrieval-oriented representations.
A criterion measuring similarity between hidden representation spaces does not itself define an optimal ordering of coordinates under truncation.
Our objective therefore does not attempt to reproduce the teacher representation itself; it attempts to reproduce the teacher's semantic neighborhood distribution at each storage rate.

\subsection{Representation Geometry}

Normalized contrastive representations naturally live on a hypersphere.
Prior work has shown that alignment and uniformity provide useful geometric perspectives on contrastive representation quality~\cite{wang2020alignment}.
Centered Kernel Alignment (CKA) provides a principled measure for comparing representational similarity structures across feature spaces of different dimensionalities~\cite{kornblith2019similarity}.

These tools address different questions from ours.
CKA asks whether two representation spaces induce similar relational structure.
SDR-MRL instead defines a local probabilistic geometry:
for each anchor, similarities induce a Gibbs distribution over semantic neighbors.
Matching these distributions focuses training on relative local structure directly relevant to nearest-neighbor retrieval.

\subsection{Information-Theoretic Representation Learning}

The information bottleneck framework separates predictive or task-relevant information from information about the raw input.
Variational approximations make such objectives tractable for neural representation learning~\cite{alemi2017vib}.
MCR$^2$ provides another information-theoretic view by optimizing finite-sample coding-rate reduction to learn discriminative low-dimensional structure~\cite{yu2020mcr2}.

Our proposal does not directly estimate $I(X;Z)$ or impose a Shannon coding constraint.
Instead, the deployment rate is operational:
if each coordinate is stored using $b$ bits, a $d_k$-dimensional embedding requires
\begin{equation}
    R_k = b d_k
\end{equation}
bits before entropy coding.
We therefore distinguish \emph{embedding storage rate} from mutual-information rate.

\subsection{Successive Refinement}

Classical successive refinement studies whether a source can first be represented coarsely and then progressively refined while remaining rate-distortion optimal at intermediate stages~\cite{equitz1991successive}.
This perspective is conceptually aligned with Matryoshka representations:
a low-dimensional prefix acts as a coarse description, while later coordinates refine the description.

However, successful refinement does not generally imply that every later block must provide less information than every earlier block.
Semantic variables may exhibit synergy.
Accordingly, our core method does \emph{not} impose a universal diminishing-marginal-information constraint.
Instead, it optimizes low distortion at all requested rates and optionally enforces only monotonic refinement.


% =========================================================
\section{Problem Formulation}
% =========================================================

Let
\begin{equation}
    f_\theta:\mathcal{X}\rightarrow\R^D
\end{equation}
be an encoder and
\begin{equation}
    h_i = f_\theta(x_i)
\end{equation}
the raw representation for sample $x_i$.

We define a sequence of deployment dimensions
\begin{equation}
    0=d_0 < d_1 < d_2 < \cdots < d_K=D.
\end{equation}

The normalized embedding at level $k$ is
\begin{equation}
    z_i^{(k)}
    =
    \frac{
        h_{i,1:d_k}
    }{
        \|h_{i,1:d_k}\|_2
    }.
    \label{eq:prefix}
\end{equation}

We use $Z_k$ for the corresponding random variable.

Although every prefix is normalized independently, the representation remains informationally nested.
In particular,
\begin{equation}
    Z_{k-1}=\psi_k(Z_k),
\end{equation}
where $\psi_k$ truncates the first $d_{k-1}$ coordinates and renormalizes them.
Therefore
\begin{equation}
    X \rightarrow Z_K \rightarrow Z_{K-1}
    \rightarrow \cdots \rightarrow Z_1
\end{equation}
forms a deterministic hierarchy.

Let
\begin{equation}
    B_k=h_{d_{k-1}+1:d_k}
\end{equation}
denote the newly available coordinate block at level $k$.
Before normalization, the nested representation can be written as
\begin{equation}
    H_k=(H_{k-1},B_k).
\end{equation}

For fixed numerical precision $b$, we define the operational storage rate
\begin{equation}
    R_k = b d_k.
    \label{eq:storage-rate}
\end{equation}
When $b$ is constant across all models, $d_k$ itself is sufficient for rate comparisons.


% =========================================================
\section{Method}
% =========================================================

\subsection{Overview}

SDR-MRL contains three conceptual components:

\begin{enumerate}[leftmargin=1.5em]
    \item a full-dimensional representation defines a semantic neighborhood distribution;
    \item every truncated prefix attempts to reconstruct this distribution;
    \item training minimizes the resulting semantic distortion over a distribution of deployment rates.
\end{enumerate}

The main training objective is intentionally small:
\begin{equation}
\boxed{
    \cL_{\mathrm{SDR}}
    =
    \cL_{\mathrm{task}}
    +
    \lambda_{\mathrm{sem}}\cL_{\mathrm{sem}}
    +
    \lambda_{\mathrm{mono}}\cL_{\mathrm{mono}}.
}
\label{eq:main-objective}
\end{equation}

The semantic term is the principal proposed component.
The monotonic term is optional and will be validated by ablation.


% ---------------------------------------------------------
\subsection{Semantic Neighborhood Teacher}
% ---------------------------------------------------------

For an anchor $x_i$, consider a candidate/reference set
\begin{equation}
    \mathcal{C}_i=\{x_j:j\neq i\}.
\end{equation}
In the simplest implementation, $\mathcal{C}_i$ contains all other examples in the minibatch.
A memory queue or retrieval corpus can be substituted when larger candidate sets are desirable.

The full representation is treated as a stop-gradient teacher:
\begin{equation}
    \bar z_i^{(K)} = \sg(z_i^{(K)}).
\end{equation}

We define teacher logits
\begin{equation}
    a_{ij}^{T}
    =
    \frac{
        \simcos(
            \bar z_i^{(K)},
            \bar z_j^{(K)}
        )
    }{\tau_T},
\end{equation}
and the semantic neighborhood distribution
\begin{equation}
\boxed{
    p_T(S=j|X=x_i,\mathcal C_i)
    =
    \frac{\exp(a_{ij}^{T})}
    {
        \sum_{r\in\mathcal C_i}
        \exp(a_{ir}^{T})
    }.
}
\label{eq:teacher}
\end{equation}

Here, $S$ is a discrete semantic random variable indicating which reference item is selected according to the teacher-induced semantic neighborhood.

For readability, we suppress explicit conditioning on $\mathcal C_i$ in the theoretical derivations below.
All information-theoretic quantities may be understood as conditioned on the candidate set.


% ---------------------------------------------------------
\subsection{Prefix Semantic Decoder}
% ---------------------------------------------------------

At prefix $k$, student logits are
\begin{equation}
    a_{ij}^{(k)}
    =
    \frac{
        \simcos(
            z_i^{(k)},
            z_j^{(k)}
        )
    }{\tau_k},
\end{equation}
with induced distribution
\begin{equation}
\boxed{
    q_k(S=j|Z_k=z_i^{(k)})
    =
    \frac{\exp(a_{ij}^{(k)})}
    {
        \sum_{r\in\mathcal C_i}
        \exp(a_{ir}^{(k)})
    }.
}
\label{eq:student}
\end{equation}

This decoder contains no additional inference-time parameters.
It interprets the local cosine geometry of a prefix as a probability distribution over semantic neighbors.


% ---------------------------------------------------------
\subsection{Semantic Neighborhood Distortion}
% ---------------------------------------------------------

We define the distortion at rate $R_k$ as
\begin{equation}
\boxed{
    \cD_k
    =
    \E_X
    \KL
    \left(
        p_T(S|X)
        \Vert
        q_k(S|Z_k)
    \right).
}
\label{eq:snd}
\end{equation}

The minibatch estimator is
\begin{equation}
    \widehat{\cD}_k
    =
    \frac{1}{B}
    \sum_{i=1}^{B}
    \sum_{j\in\mathcal C_i}
    p_{ij}^{T}
    \log
    \frac{
        p_{ij}^{T}
    }{
        q_{ij}^{(k)}
    }.
\end{equation}

Forward KL is the default because it penalizes a compressed prefix for assigning insufficient mass to semantic neighbors that the teacher considers important.
Reverse KL and Jensen--Shannon divergence will be included as ablations.


% ---------------------------------------------------------
\subsection{Variational Information Interpretation}
% ---------------------------------------------------------

The teacher distribution defines the joint distribution
\begin{equation}
    p_T(x,s)=p(x)p_T(s|x).
\end{equation}

Because $Z_k$ is a deterministic function of $X$, we have the Markov structure
\begin{equation}
    S-X-Z_k.
\end{equation}

\begin{proposition}[Variational semantic distortion]
For any semantic decoder $q_k(S|Z_k)$,
\begin{equation}
\boxed{
    \cD_k
    =
    I_T(S;X|Z_k)
    +
    \epsilon_k,
}
\label{eq:variational-decomposition}
\end{equation}
where
\begin{equation}
    \epsilon_k
    =
    \E_{Z_k}
    \KL
    \left(
        p_T(S|Z_k)
        \Vert
        q_k(S|Z_k)
    \right)
    \geq 0.
\end{equation}
Consequently,
\begin{equation}
    \cD_k \geq I_T(S;X|Z_k).
\end{equation}
\end{proposition}

\begin{proof}
Starting from Eq.~\eqref{eq:snd},
\begin{align}
\cD_k
&=
\E
\log
\frac{
    p_T(S|X)
}{
    q_k(S|Z_k)
}
\\
&=
\E
\log
\frac{
    p_T(S|X)
}{
    p_T(S|Z_k)
}
+
\E
\log
\frac{
    p_T(S|Z_k)
}{
    q_k(S|Z_k)
}.
\end{align}
Because $Z_k$ is determined by $X$,
\begin{equation}
    p_T(S|X,Z_k)=p_T(S|X),
\end{equation}
and therefore the first term is
\begin{equation}
    I_T(S;X|Z_k).
\end{equation}
The second term is
\begin{equation}
    \E_{Z_k}
    \KL
    \left(
        p_T(S|Z_k)
        \Vert
        q_k(S|Z_k)
    \right).
\end{equation}
\end{proof}

This proposition gives Eq.~\eqref{eq:snd} an information-theoretic interpretation without requiring an explicit mutual-information estimator.
The practical cosine decoder produces a variational upper bound on information loss rather than a direct estimate of high-dimensional mutual information.


% ---------------------------------------------------------
\subsection{Optimal Semantic Distortion}
% ---------------------------------------------------------

Define the optimal distortion associated with representation $Z_k$ as
\begin{equation}
    \cD_k^\star
    =
    \inf_{q_k}
    \E
    \KL
    \left(
        p_T(S|X)
        \Vert
        q_k(S|Z_k)
    \right).
\end{equation}

\begin{corollary}
The optimal semantic distortion is
\begin{equation}
\boxed{
    \cD_k^\star
    =
    I_T(S;X|Z_k).
}
\label{eq:optimal-distortion}
\end{equation}
\end{corollary}

The minimizing decoder is
\begin{equation}
    q_k^\star(S|Z_k)=p_T(S|Z_k).
\end{equation}


% ---------------------------------------------------------
\subsection{Successive Semantic Refinement}
% ---------------------------------------------------------

Because the Matryoshka hierarchy is nested,
\begin{equation}
    Z_{k-1}=\psi_k(Z_k).
\end{equation}

Hence increasing the representation rate cannot increase the \emph{optimal} conditional semantic information loss.

\begin{proposition}[Monotonic optimal refinement]
For adjacent Matryoshka levels,
\begin{equation}
\boxed{
    \cD_k^\star
    \leq
    \cD_{k-1}^\star.
}
\end{equation}
Furthermore,
\begin{equation}
\boxed{
    \cD_{k-1}^\star-\cD_k^\star
    =
    I_T(S;Z_k|Z_{k-1}).
}
\label{eq:refinement-gain}
\end{equation}
\end{proposition}

If the newly available information is represented as block $B_k$ such that
\begin{equation}
    Z_k=(Z_{k-1},B_k),
\end{equation}
then
\begin{equation}
\boxed{
    \cD_{k-1}^\star-\cD_k^\star
    =
    I_T(S;B_k|Z_{k-1}).
}
\label{eq:block-information}
\end{equation}

Equation~\eqref{eq:block-information} is the key successive-refinement interpretation:
the reduction in semantic distortion after increasing the embedding rate equals the new semantic information contributed by the added representation block, conditioned on information already available at the lower rate.

For the practical decoder,
\begin{equation}
    \cD_{k-1}-\cD_k
    =
    I_T(S;B_k|Z_{k-1})
    +
    (\epsilon_{k-1}-\epsilon_k).
\end{equation}
Therefore practical distortion reductions should be treated as \emph{surrogates} for conditional semantic gains rather than exact mutual-information estimates.


% ---------------------------------------------------------
\subsection{Multi-Rate Semantic Objective}
% ---------------------------------------------------------

Let
\begin{equation}
    \pi_k
    =
    P(R=R_k)
\end{equation}
represent the probability that deployment requests rate $R_k$.

We minimize expected semantic distortion:
\begin{equation}
\boxed{
    \cL_{\mathrm{sem}}
    =
    \sum_{k=1}^{K-1}
    \pi_k \cD_k.
}
\label{eq:semantic-objective}
\end{equation}

This formulation gives prefix weights an operational interpretation.
If the deployment distribution is unknown, we use
\begin{equation}
    \pi_k=\frac{1}{K-1}.
\end{equation}
For low-rate-focused systems, alternatives such as
\begin{equation}
    \pi_k \propto d_k^{-1}
\end{equation}
will be investigated.


% ---------------------------------------------------------
\subsection{Task Objective}
% ---------------------------------------------------------

The semantic self-distillation term alone admits degenerate solutions if the teacher representation itself collapses.
We therefore retain an ordinary task objective.

For paired retrieval training, one possible prefix-specific loss is
\begin{equation}
\cL_{\mathrm{ret}}^{(k)}
=
-\frac{1}{B}
\sum_{i=1}^{B}
\log
\frac{
    \exp(
        \simcos(
            z_i^{(k)},
            z_{i^+}^{(k)}
        )/\tau_R
    )
}{
    \sum_{j\in\mathcal C_i}
    \exp(
        \simcos(
            z_i^{(k)},
            z_j^{(k)}
        )/\tau_R
    )
}.
\end{equation}

A standard MRL task objective is then
\begin{equation}
    \cL_{\mathrm{task}}
    =
    \sum_{k=1}^{K}
    \alpha_k
    \cL_{\mathrm{task}}^{(k)}.
\end{equation}

Importantly, when a contrastive task loss is used here, it functions as downstream retrieval supervision; we do not interpret its InfoNCE form as a quantitative estimator of hidden-state mutual information.


% ---------------------------------------------------------
\subsection{Monotonic Refinement Regularization}
% ---------------------------------------------------------

Although optimal distortion must be monotone in the nested representation rate, the practical cosine decoder can violate this property:
\begin{equation}
    \cD_k > \cD_{k-1}.
\end{equation}

We therefore consider the optional regularizer
\begin{equation}
\boxed{
    \cL_{\mathrm{mono}}
    =
    \sum_{k=2}^{K-1}
    \left[
        \cD_k-\cD_{k-1}
    \right]_+.
}
\label{eq:mono}
\end{equation}

This term discourages cases where adding coordinates makes semantic neighborhood recovery worse.

We deliberately do \emph{not} require marginal gains to decrease with $k$.
Such a condition holds in the linear-Gaussian analysis below but is not universal.
For example, an XOR semantic variable may contain purely synergistic information for which a later block is more informative conditional on an earlier block than either block is independently.


% ---------------------------------------------------------
\subsection{Final Training Objective}
% ---------------------------------------------------------

The complete objective is
\begin{equation}
\boxed{
    \cL_{\mathrm{SDR}}
    =
    \cL_{\mathrm{task}}
    +
    \lambda_{\mathrm{sem}}
    \sum_{k=1}^{K-1}
    \pi_k\cD_k
    +
    \lambda_{\mathrm{mono}}
    \sum_{k=2}^{K-1}
    [\cD_k-\cD_{k-1}]_+.
}
\end{equation}

The minimal proposed model sets
\begin{equation}
    \lambda_{\mathrm{mono}}=0
\end{equation}
and tests semantic neighborhood distortion alone.
The monotonic regularizer is retained only if its ablation consistently improves the distortion-rate profile or downstream low-rate performance.


% ---------------------------------------------------------
\subsection{Geometric Interpretation of the Gradient}
% ---------------------------------------------------------

For an anchor $i$, let
\begin{equation}
    p_{ij}=p_T(S=j|X=x_i),
    \qquad
    q_{ij}=q_k(S=j|Z_k=z_i).
\end{equation}

Because
\begin{equation}
    \cD_i
    =
    \KL(p_i\Vert q_i),
\end{equation}
the derivative with respect to a student logit is
\begin{equation}
\boxed{
    \frac{\partial\cD_i}
    {\partial a_{ij}^{(k)}}
    =
    q_{ij}-p_{ij}.
}
\end{equation}

For normalized embeddings,
\begin{equation}
    a_{ij}^{(k)}
    =
    \frac{
        (z_i^{(k)})^\top z_j^{(k)}
    }{\tau_k},
\end{equation}
yielding the direct anchor gradient
\begin{equation}
\boxed{
    \frac{\partial\cD_i}
    {\partial z_i^{(k)}}
    =
    \frac{1}{\tau_k}
    \sum_j
    (q_{ij}-p_{ij})z_j^{(k)}.
}
\label{eq:geometry-gradient}
\end{equation}

The gradient therefore moves the student's local semantic barycenter
\begin{equation}
    \sum_j q_{ij}z_j
\end{equation}
toward the teacher-weighted semantic barycenter
\begin{equation}
    \sum_j p_{ij}z_j.
\end{equation}

For raw prefix $h_i$ and
\begin{equation}
    z_i=\frac{h_i}{\|h_i\|},
\end{equation}
the gradient is projected onto the tangent space of the unit hypersphere:
\begin{equation}
\boxed{
    \frac{\partial\cD_i}{\partial h_i}
    =
    \frac{1}{
        \tau_k\|h_i\|
    }
    (I-z_i z_i^\top)
    \sum_j(q_{ij}-p_{ij})z_j.
}
\end{equation}

Thus the proposed semantic loss performs a local geometric correction directly on the normalized embedding manifold.


% ---------------------------------------------------------
\subsection{Efficient Stochastic Rate Sampling}
% ---------------------------------------------------------

Computing semantic distortion at every prefix is unnecessary.

If
\begin{equation}
    K_t\sim\pi,
\end{equation}
then
\begin{equation}
    \widehat{\cL}_{\mathrm{sem}}
    =
    \cD_{K_t}
\end{equation}
is an unbiased estimator of Eq.~\eqref{eq:semantic-objective}:
\begin{equation}
    \E_{K_t\sim\pi}
    [
        \cD_{K_t}
    ]
    =
    \sum_k\pi_k\cD_k.
\end{equation}

For monotonic regularization, we analogously sample an adjacent edge
\begin{equation}
    (k-1,k)
\end{equation}
and compute only the two associated prefix distortions.

This stochastic variant reduces training cost while retaining the desired multi-rate objective in expectation.


% ---------------------------------------------------------
\subsection{Training Algorithm}
% ---------------------------------------------------------

\begin{algorithm}[t]
\caption{SDR-MRL training with stochastic rate sampling}
\label{alg:sdr}
\begin{algorithmic}[1]

\Require batch $\mathcal B$, encoder $f_\theta$, dimensions $\{d_k\}$,
rate distribution $\pi$, temperatures $\tau_T,\tau_S$

\State $H \gets f_\theta(\mathcal B)$
\State $Z_K \gets \operatorname{Normalize}(H_{:,1:D})$

\State $P_T \gets
\operatorname{StopGrad}
\left[
\operatorname{Softmax}
(
Z_K Z_K^\top/\tau_T
)
\right]$

\State sample $k\sim\pi$

\State $Z_k \gets
\operatorname{Normalize}
(
H_{:,1:d_k}
)$

\State $Q_k \gets
\operatorname{Softmax}
(
Z_k Z_k^\top/\tau_S
)$

\State $\cD_k \gets \KL(P_T\Vert Q_k)$

\State compute $\cL_{\mathrm{task}}$

\If{$\lambda_{\mathrm{mono}}>0$ and $k>1$}
    \State construct $Z_{k-1}$ and $Q_{k-1}$
    \State $\cD_{k-1}\gets\KL(P_T\Vert Q_{k-1})$
    \State $\cL_{\mathrm{mono}}
    \gets
    [\cD_k-\cD_{k-1}]_+$
\Else
    \State $\cL_{\mathrm{mono}}\gets0$
\EndIf

\State $\cL
\gets
\cL_{\mathrm{task}}
+
\lambda_{\mathrm{sem}}\cD_k
+
\lambda_{\mathrm{mono}}\cL_{\mathrm{mono}}$

\State update $\theta$ using $\nabla_\theta\cL$

\end{algorithmic}
\end{algorithm}


% ---------------------------------------------------------
\subsection{Teacher Stability}
% ---------------------------------------------------------

A self-teacher can drift during training.
Our default configuration uses a stop-gradient full-dimensional teacher together with a full-dimensional task loss.
We will additionally test:

\begin{enumerate}[leftmargin=1.5em]
    \item \textbf{Online self-teacher:}
    current full-dimensional representation with stop-gradient;

    \item \textbf{EMA teacher:}
    a slowly updated copy of the encoder;

    \item \textbf{Frozen teacher:}
    an independently trained full-dimensional model.
\end{enumerate}

The online teacher is preferred if performance is comparable because it adds no persistent model at deployment or training initialization.


% =========================================================
\section{Linear-Gaussian Analysis}
% =========================================================

The variational results above apply to arbitrary nested representations.
We next study a tractable setting in which the optimal semantic representation has an explicit Matryoshka ordering.

Let
\begin{equation}
    X\in\R^p,
    \qquad
    S\in\R^q
\end{equation}
be jointly Gaussian.
Assume
\begin{equation}
    \Sigma_X\succ0
\end{equation}
and define the whitened input
\begin{equation}
    \widetilde X
    =
    \Sigma_X^{-1/2}X,
    \qquad
    \operatorname{Cov}(\widetilde X)=I.
\end{equation}

Consider rank-$d$ linear representations
\begin{equation}
    Z_d=V_d^\top\widetilde X,
    \qquad
    V_d^\top V_d=I_d.
\end{equation}

We evaluate semantic squared reconstruction distortion:
\begin{equation}
    D(V_d)
    =
    \E
    \left[
        \|
            S-\E[S|Z_d]
        \|_2^2
    \right].
\end{equation}

For jointly Gaussian variables, the MMSE decoder is linear and
\begin{equation}
\begin{aligned}
    D(V_d)
    &=
    \operatorname{tr}(\Sigma_S)
    \\
    &\quad -
    \operatorname{tr}
    \left(
        V_d^\top
        M
        V_d
    \right),
\end{aligned}
\end{equation}
where
\begin{equation}
\boxed{
    M
    =
    \Sigma_X^{-1/2}
    \Sigma_{XS}
    \Sigma_{SX}
    \Sigma_X^{-1/2}.
}
\label{eq:predictive-matrix}
\end{equation}

\begin{theorem}[Nested semantic optimum]
Let
\begin{equation}
    \lambda_1\geq
    \lambda_2\geq\cdots\geq0
\end{equation}
be the eigenvalues of $M$, with corresponding orthonormal eigenvectors
\begin{equation}
    u_1,u_2,\ldots.
\end{equation}
Then an optimal rank-$d$ semantic representation is obtained with
\begin{equation}
\boxed{
    V_d^\star
    =
    [u_1,\ldots,u_d].
}
\end{equation}

Consequently,
\begin{equation}
    \operatorname{span}(V_1^\star)
    \subset
    \operatorname{span}(V_2^\star)
    \subset
    \cdots,
\end{equation}
so the family of semantic-optimal linear representations admits a Matryoshka ordering.
\end{theorem}

\begin{proof}
Minimizing $D(V_d)$ is equivalent to
\begin{equation}
    \max_{V_d^\top V_d=I}
    \operatorname{tr}
    (
        V_d^\top M V_d
    ).
\end{equation}
By the Ky Fan maximum principle, this trace is maximized by the top-$d$ eigenvectors of $M$.
\end{proof}

The optimal distortion is
\begin{equation}
\boxed{
    D_d^\star
    =
    \operatorname{tr}(\Sigma_S)
    -
    \sum_{j=1}^d\lambda_j.
}
\end{equation}

Therefore the semantic gain of the $d$th coordinate is
\begin{equation}
\boxed{
    D_{d-1}^\star-D_d^\star
    =
    \lambda_d.
}
\end{equation}

This provides a theoretical example in which semantic information is naturally ordered from coarse to fine.

Importantly, the relevant directions depend on
\begin{equation}
    \Sigma_{XS}\Sigma_{SX}
\end{equation}
rather than input variance alone.
Hence variance-maximizing directions such as PCA are optimal only in special cases.
A high-variance nuisance direction can be less useful than a lower-variance direction that is more predictive of $S$.

\paragraph{Scope of the theorem.}
The decreasing eigenvalue sequence does not justify a universal diminishing-gain constraint for arbitrary semantic representations.
Nonlinear semantic variables can contain synergistic information.
The theorem therefore motivates semantic ordering while the general method requires only low distortion across nested rates.


% =========================================================
\section{Semantic Distortion-Rate Evaluation}
% =========================================================

\subsection{Distortion-Rate Profile}

Given storage rates
\begin{equation}
    R_k=bd_k
\end{equation}
and a common semantic reference teacher $T_{\mathrm{ref}}$, define
\begin{equation}
    \cD_k^{\mathrm{ref}}
    =
    \E
    \KL
    \left(
        p_{T_{\mathrm{ref}}}(S|X)
        \Vert
        q_k(S|Z_k)
    \right).
\end{equation}

The semantic distortion-rate profile is
\begin{equation}
\boxed{
    \mathcal P
    =
    \{
        (R_k,\cD_k^{\mathrm{ref}})
    \}_{k=1}^{K}.
}
\end{equation}

The reference teacher is held fixed across all compared models.
This prevents a weak model from appearing to have low distortion merely because its own prefixes agree with its own weak full-dimensional representation.


% ---------------------------------------------------------
\subsection{Normalized Distortion Area}
% ---------------------------------------------------------

Let $q_0$ denote a zero-rate uniform semantic decoder.
Define
\begin{equation}
    \cD_0^{\mathrm{ref}}
    =
    \E
    \KL(
        p_{T_{\mathrm{ref}}}
        \Vert q_0
    ).
\end{equation}

A normalized distortion can be written as
\begin{equation}
    \widetilde{\cD}_k
    =
    \frac{
        \cD_k^{\mathrm{ref}}
        -
        \cD_D^{\mathrm{ref}}
    }{
        \cD_0^{\mathrm{ref}}
        -
        \cD_D^{\mathrm{ref}}
        +\varepsilon
    }.
\end{equation}

With normalized rate
\begin{equation}
    r_k=d_k/D,
\end{equation}
we compute the semantic distortion-rate area
\begin{equation}
\boxed{
    \mathrm{SDRA}
    =
    \sum_{k=1}^{K}
    \frac{
        \widetilde{\cD}_{k-1}
        +
        \widetilde{\cD}_k
    }{2}
    (r_k-r_{k-1}).
}
\end{equation}

Lower SDRA indicates that useful semantic structure becomes recoverable earlier in the representation.


% ---------------------------------------------------------
\subsection{Price of Nestedness}
% ---------------------------------------------------------

Let $f_d^\star$ denote a model independently optimized at dimension $d$ and
$f_{\mathrm{nested}}^{1:d}$ the $d$-dimensional prefix of a single Matryoshka model.

For downstream quality metric $Q$,
\begin{equation}
\boxed{
    P_{\mathrm{nest}}(d)
    =
    Q(f_d^\star)
    -
    Q(f_{\mathrm{nested}}^{1:d}).
}
\end{equation}

Aggregate price of nestedness is
\begin{equation}
    P_{\mathrm{nest}}
    =
    \sum_k
    \pi_k P_{\mathrm{nest}}(d_k).
\end{equation}

This directly measures the performance paid for requiring all deployment capacities to coexist in one representation.


% ---------------------------------------------------------
\subsection{Neighborhood Preservation}
% ---------------------------------------------------------

For reference teacher neighborhood
\begin{equation}
    \mathcal N_M^T(i)
\end{equation}
and prefix neighborhood
\begin{equation}
    \mathcal N_M^{(k)}(i),
\end{equation}
we measure
\begin{equation}
\boxed{
\mathrm{kNNRecall}_k
=
\frac{1}{N}
\sum_{i=1}^N
\frac{
    |
    \mathcal N_M^T(i)
    \cap
    \mathcal N_M^{(k)}(i)
    |
}{
    M
}.
}
\end{equation}

We will additionally report:
\begin{itemize}[leftmargin=1.5em]
    \item Spearman correlation of pairwise similarities,
    \item neighborhood trustworthiness,
    \item neighborhood continuity,
    \item downstream retrieval ranking metrics.
\end{itemize}


% =========================================================
\section{Experimental Plan}
% =========================================================

\subsection{Research Questions}

The experiments will test the following questions:

\begin{description}[leftmargin=1.6em,style=nextline]

\item[RQ1.]
Does semantic neighborhood distortion improve extreme low-dimensional Matryoshka embeddings beyond standard MRL?

\item[RQ2.]
Does neighborhood-distribution matching outperform generic representation alignment such as pairwise similarity matching or CKA?

\item[RQ3.]
Does adding monotonic refinement improve the semantic distortion-rate profile?

\item[RQ4.]
Does semantic distortion reduction
\begin{equation}
    \cD_{k-1}-\cD_k
\end{equation}
predict actual downstream gains from increasing representation dimensionality?

\item[RQ5.]
Can stochastic rate sampling preserve performance while lowering training overhead relative to full multi-prefix alignment?

\item[RQ6.]
Does SDR-MRL reduce the price of nestedness relative to MRL and existing Matryoshka compression methods?

\end{description}


% ---------------------------------------------------------
\subsection{Backbones}
% ---------------------------------------------------------

We propose three experimental scales.

\paragraph{Controlled setting.}
BERT-base or an equivalent medium-sized sentence encoder will be used for rapid iteration and direct low-cost ablation.

\paragraph{Retrieval-oriented setting.}
BGE-M3 will serve as a modern embedding backbone because it has been evaluated in recent Matryoshka work and supports retrieval-oriented experiments.

\paragraph{Scaling setting.}
A compact Qwen3-family backbone can be included after the method is stabilized to test whether the behavior persists in a larger contemporary model.

The initial paper does not require all three scales if compute is constrained.
The controlled and retrieval-oriented settings are sufficient for the primary hypothesis.


% ---------------------------------------------------------
\subsection{Prefix Dimensions}
% ---------------------------------------------------------

For a $D=768$ backbone:
\begin{equation}
    \mathcal D
    =
    \{
        16,32,64,128,256,512,768
    \}.
\end{equation}

For architectures with different full dimensionality, the same approximate geometric progression will be used.

The primary regime of interest is
\begin{equation}
    d\in\{16,32,64\},
\end{equation}
where representational bottlenecks are most severe.


% ---------------------------------------------------------
\subsection{Benchmarks}
% ---------------------------------------------------------

\paragraph{Semantic textual similarity.}
We will evaluate STS12--STS16, STS-B, and SICK-R to maintain comparability with sentence-embedding literature.

\paragraph{Retrieval.}
The central evaluation will use BEIR and selected MTEB retrieval tasks.
Metrics will include
\begin{equation}
    \mathrm{nDCG}@10,
    \qquad
    \mathrm{Recall}@k.
\end{equation}

\paragraph{Classification and NLI.}
A smaller classification/NLI suite will be retained for direct comparison to recent Matryoshka methods, but retrieval is the primary target because neighborhood geometry is the central object of the proposed method.

\paragraph{Out-of-domain evaluation.}
Where possible, the reference teacher will be trained on the source training distribution while semantic distortion and retrieval quality are evaluated under domain shift.


% ---------------------------------------------------------
\subsection{Baselines}
% ---------------------------------------------------------

The main comparison set is:

\begin{enumerate}[leftmargin=1.5em]
    \item independently trained fixed-dimensional embeddings;
    \item vanilla MRL~\cite{kusupati2022matryoshka};
    \item ESE~\cite{li2025ese};
    \item Matryoshka-Adaptor~\cite{yoon2024matryoshka};
    \item SMEC~\cite{zhang2025smec};
    \item MIPIC~\cite{huy2026mipic};
    \item SDR-MRL.
\end{enumerate}

Not every baseline is applicable to every backbone or supervision setting.
Comparisons will therefore be reported only where training assumptions are compatible.


% ---------------------------------------------------------
\subsection{Primary Metrics}
% ---------------------------------------------------------

For every representation dimension, we will report:

\begin{itemize}[leftmargin=1.5em]
    \item downstream task quality;
    \item fixed-teacher semantic distortion $\cD_k^{\mathrm{ref}}$;
    \item kNN neighborhood recall;
    \item embedding storage;
    \item retrieval latency where applicable;
    \item training throughput;
    \item price of nestedness.
\end{itemize}

Aggregate metrics include:
\begin{itemize}[leftmargin=1.5em]
    \item average low-rate task performance over $16/32/64$ dimensions;
    \item semantic distortion-rate area;
    \item deployment-prior-weighted performance;
    \item deployment-prior-weighted price of nestedness.
\end{itemize}


% =========================================================
\section{Ablation Study}
% =========================================================

Ablation is central to the proposal because several alternative explanations must be separated:
the gain may come from ordinary self-distillation, from geometric neighborhood structure, from the information-theoretic formulation, or simply from additional training compute.

\begin{table*}[t]
\centering
\small
\begin{tabularx}{\textwidth}{
    >{\raggedright\arraybackslash}p{0.07\textwidth}
    >{\raggedright\arraybackslash}p{0.22\textwidth}
    >{\raggedright\arraybackslash}X
    >{\raggedright\arraybackslash}X
}
\toprule
\textbf{ID}
&
\textbf{Variant}
&
\textbf{Question}
&
\textbf{Interpretation}
\\
\midrule

A0 &
MRL only &
What is the nested baseline? &
Establishes ordinary multi-prefix supervision.
\\

A1 &
MRL + SND &
Does semantic neighborhood distortion alone help? &
Core test of the proposal.
\\

A2 &
A1 + monotonic refinement &
Does enforcing $\cD_k\leq\cD_{k-1}$ help practical prefixes? &
Tests whether decoder-level non-monotonicity is a meaningful failure mode.
\\

A3 &
Self vs EMA vs frozen teacher &
Is the result dependent on teacher drift? &
Separates online self-distillation effects from semantic-geometry preservation.
\\

A4 &
Forward KL / reverse KL / JS &
Which probabilistic geometry is appropriate? &
Tests mass-covering versus mode-seeking neighborhood preservation.
\\

A5 &
SND vs Gram-MSE vs CKA &
Is distributional semantic geometry better than generic representation alignment? &
Critical comparison to alternative geometric objectives.
\\

A6 &
All candidates / teacher top-$M$ / union hard neighbors &
How local should the semantic graph be? &
Measures accuracy--cost trade-off and false-neighbor sensitivity.
\\

A7 &
Uniform $\pi$ / inverse-dimension $\pi$ / deployment prior &
Does rate weighting control information allocation as predicted? &
Tests operational interpretation of prefix weights.
\\

A8 &
All-prefix training vs sampled-rate training &
Can the objective be optimized efficiently? &
Measures performance versus training throughput.
\\

A9 &
Temperature sweep &
How sensitive is the semantic distribution to teacher entropy? &
Checks whether gains depend on a narrowly tuned neighborhood sharpness.
\\

A10 &
Fixed full-dimensional teacher quality &
Does a stronger semantic teacher improve compressed prefixes? &
Tests whether full-space semantic quality is the limiting factor.
\\

A11 &
Independent low-dimensional oracle &
How large is the price of nestedness? &
Measures how close the nested representation approaches separately optimized models.
\\

\bottomrule
\end{tabularx}
\caption{
Planned ablations for SDR-MRL.
}
\label{tab:ablations}
\end{table*}


% ---------------------------------------------------------
\subsection{Ablation 1: Is Semantic Neighborhood Distortion Sufficient?}
% ---------------------------------------------------------

The first experiment should compare only
\begin{equation}
    \cL_{\mathrm{MRL}}
\end{equation}
against
\begin{equation}
    \cL_{\mathrm{MRL}}
    +
    \lambda_{\mathrm{sem}}\cL_{\mathrm{sem}}.
\end{equation}

This is the principal go/no-go experiment.

The research direction should be considered promising only if SND produces a consistent improvement at low dimensions, especially
\begin{equation}
    d\in\{16,32,64\},
\end{equation}
without materially degrading the full-dimensional representation.


% ---------------------------------------------------------
\subsection{Ablation 2: Geometry Objective}
% ---------------------------------------------------------

We compare four forms of cross-dimensional structural preservation.

\paragraph{Semantic neighborhood KL.}
\begin{equation}
    \KL(P_T\Vert Q_k).
\end{equation}

\paragraph{Pairwise similarity matching.}
\begin{equation}
    \|
        G_T-G_k
    \|_F^2,
\end{equation}
where $G$ is a cosine Gram matrix.

\paragraph{CKA.}
Linear CKA between full and truncated representations.

\paragraph{Hard-neighbor supervision.}
Cross entropy using only the teacher's nearest neighbor.

If SND improves both retrieval and independently measured neighborhood preservation beyond these alternatives, it supports the claim that probabilistic local semantic geometry is useful beyond generic feature alignment.


% ---------------------------------------------------------
\subsection{Ablation 3: Semantic Teacher}
% ---------------------------------------------------------

Three teachers will be compared:

\begin{equation}
    T_{\mathrm{online}},
    \qquad
    T_{\mathrm{EMA}},
    \qquad
    T_{\mathrm{frozen}}.
\end{equation}

The following should be measured:

\begin{itemize}[leftmargin=1.5em]
    \item full-dimensional retrieval quality;
    \item average teacher neighborhood entropy;
    \item student distortion;
    \item low-dimensional downstream quality.
\end{itemize}

A large difference between online and frozen teachers would indicate that teacher drift is a major confound and should be treated explicitly in the final method.


% ---------------------------------------------------------
\subsection{Ablation 4: Monotonic Refinement}
% ---------------------------------------------------------

We test
\begin{equation}
    \lambda_{\mathrm{mono}}
    \in
    \{0,0.05,0.1,0.3\}.
\end{equation}

Besides downstream metrics, we measure the monotonicity violation rate
\begin{equation}
\boxed{
    V_{\mathrm{mono}}
    =
    \frac{1}{K-2}
    \sum_{k=2}^{K-1}
    \mathbf{1}
    [
        \cD_k>\cD_{k-1}
    ].
}
\end{equation}

The regularizer should remain in the final method only if reducing $V_{\mathrm{mono}}$ is associated with better downstream distortion-rate curves.


% ---------------------------------------------------------
\subsection{Ablation 5: Deployment-Rate Prior}
% ---------------------------------------------------------

We compare:
\begin{align}
    \pi_k^{\mathrm{uniform}}
    &=
    \frac{1}{K-1},
    \\
    \pi_k^{\mathrm{low-rate}}
    &\propto
    \frac{1}{d_k},
\end{align}
and an application-specific deployment prior.

The prediction is that low-rate weighting should shift semantic information toward earlier prefixes and improve $16$--$64$ dimensional performance, potentially at a small cost to intermediate rates.


% ---------------------------------------------------------
\subsection{Ablation 6: Candidate Neighborhood}
% ---------------------------------------------------------

We compare:

\begin{enumerate}[leftmargin=1.5em]
    \item all in-batch candidates;
    \item teacher top-$M$ candidates;
    \item teacher top-$M$ plus student hard negatives.
\end{enumerate}

The third variant addresses a failure mode in which a compressed student assigns high similarity to an incorrect point absent from the teacher's top-$M$ support.


% ---------------------------------------------------------
\subsection{Ablation 7: Temperature}
% ---------------------------------------------------------

Proposed sweep:
\begin{equation}
    \tau_T,\tau_S
    \in
    \{
        0.03,0.05,0.10,0.20
    \}.
\end{equation}

We will report teacher neighborhood entropy
\begin{equation}
    H(P_T)
\end{equation}
to determine whether failure arises from overly sharp or overly uniform semantic distributions.


% ---------------------------------------------------------
\subsection{Ablation 8: Rate Sampling and Efficiency}
% ---------------------------------------------------------

Compare full-prefix computation to stochastic rate sampling.

Report:
\begin{equation}
    \text{iterations/sec},
    \qquad
    \text{samples/sec},
    \qquad
    \text{GPU memory},
\end{equation}
together with low-rate quality and SDRA.

A desirable outcome is performance statistically indistinguishable from full-prefix optimization at substantially lower training cost.


% =========================================================
\section{Diagnostic Experiments}
% =========================================================

\subsection{Random Rotation Stress Test}

To demonstrate that full-space geometry does not determine prefix quality, construct
\begin{equation}
    Z'=ZQ
\end{equation}
for random orthogonal matrix
\begin{equation}
    Q^\top Q=I.
\end{equation}

Full-dimensional inner products satisfy
\begin{equation}
    Z'Z'^\top=ZZ^\top.
\end{equation}

Thus full-space cosine geometry is preserved.

However,
\begin{equation}
    Z'_{1:d}
\end{equation}
is generally not equivalent to
\begin{equation}
    Z_{1:d}.
\end{equation}

We will measure low-dimensional retrieval after multiple random rotations while confirming invariant full-dimensional retrieval.
This experiment isolates the importance of \emph{coordinate information ordering} from full-space representation quality.


% ---------------------------------------------------------
\subsection{Does Distortion Reduction Predict Utility?}
% ---------------------------------------------------------

For each model, dataset, and adjacent prefix pair, define
\begin{equation}
    G_k
    =
    \cD_{k-1}^{\mathrm{ref}}
    -
    \cD_k^{\mathrm{ref}}
\end{equation}
and downstream gain
\begin{equation}
    \Delta Q_k
    =
    Q_k-Q_{k-1}.
\end{equation}

We test the rank correlation
\begin{equation}
\boxed{
    \rho
    =
    \operatorname{Spearman}
    (
        G_k,\Delta Q_k
    ).
}
\end{equation}

A strong positive correlation would support the interpretation of distortion reduction as a practical measure of incremental semantic value.


% ---------------------------------------------------------
\subsection{Semantic Information Ordering}
% ---------------------------------------------------------

For analysis only, define gain per added coordinate:
\begin{equation}
    \eta_k
    =
    \frac{
        \cD_{k-1}-\cD_k
    }{
        d_k-d_{k-1}
    }.
\end{equation}

We do not enforce monotonicity of $\eta_k$.
Instead, we plot it to understand where different models allocate semantic improvement across coordinates.

This diagnostic directly tests whether MRL, MIPIC, ESE, and SDR-MRL learn qualitatively different information-allocation patterns.


% =========================================================
\section{Implementation Plan}
% =========================================================

\subsection{Initial Configuration}

The first implementation will use:

\begin{align}
\mathcal D
&=
\{16,32,64,128,256,512,D\},
\\
\pi_k
&=
1/(K-1),
\\
\tau_T
&=
0.05,
\\
\tau_S
&=
0.05.
\end{align}

Initial hyperparameter search:
\begin{align}
\lambda_{\mathrm{sem}}
&\in
\{0.1,0.3,1.0\},
\\
\lambda_{\mathrm{mono}}
&\in
\{0,0.05,0.1,0.3\}.
\end{align}

The initial candidate set should contain all in-batch examples.
Top-$M$ semantic graphs should only be introduced if all-pairs computation becomes a bottleneck.


% ---------------------------------------------------------
\subsection{Recommended Development Order}

The experiments should be implemented in the following order:

\begin{enumerate}[leftmargin=1.5em]

    \item reproduce MRL at all target dimensions;

    \item reproduce or obtain a reliable MIPIC baseline in the controlled setting;

    \item implement semantic neighborhood distortion only;

    \item evaluate low-dimensional retrieval and fixed-teacher neighborhood preservation;

    \item add monotonic refinement only if needed;

    \item implement stochastic rate sampling;

    \item scale to retrieval-oriented backbones;

    \item perform theory-specific diagnostics and random-rotation analysis.

\end{enumerate}

This ordering prevents theoretical complexity from being developed around a method that does not first demonstrate a strong empirical signal.


% =========================================================
\section{Falsification Criteria}
% =========================================================

The proposal should be reconsidered if the following occur:

\begin{enumerate}[leftmargin=1.5em]

    \item Semantic neighborhood distortion does not consistently improve $16/32/64$ dimensional task quality relative to MRL.

    \item Improvements vanish when compute-matched against pairwise or CKA distillation.

    \item Fixed-teacher semantic distortion improves but downstream retrieval does not, indicating that the chosen teacher neighborhood is not an appropriate semantic variable.

    \item Distortion reduction has little or no relationship with downstream gains across dimensions.

    \item The method requires a very narrow teacher temperature or extremely strong external teacher to function.

    \item Training overhead becomes comparable to or larger than more complex alignment methods without a corresponding performance advantage.

\end{enumerate}

These criteria ensure that the information-theoretic interpretation remains empirically meaningful rather than merely serving as an alternative notation for self-distillation.


% =========================================================
\section{Expected Contributions and Positioning}
% =========================================================

The proposed contribution should not be positioned as simply replacing CKA or InfoNCE with another auxiliary loss.

The intended conceptual distinction is:

\begin{equation}
\boxed{
\begin{array}{c}
\text{Existing alignment view:}\\[2mm]
\text{How similar is the compressed representation}\\
\text{to a richer representation?}
\end{array}
}
\end{equation}

versus

\begin{equation}
\boxed{
\begin{array}{c}
\text{SDR-MRL view:}\\[2mm]
\text{At a given representation rate,}\\
\text{how much semantic relation remains recoverable?}
\end{array}
}
\end{equation}

MIPIC provides a particularly strong comparison point because it explicitly improves cross-dimensional geometric consistency and depth-wise information flow.
SDR-MRL instead focuses on \emph{rate-conditioned semantic information}:
\begin{equation}
    I_T(S;X|Z_k)
\end{equation}
and
\begin{equation}
    I_T(S;B_k|Z_{k-1}).
\end{equation}

The desired final paper claim is therefore:

\begin{quote}
\emph{
Matryoshka learning should be understood not only as multi-prefix supervision or representation alignment, but as the construction of a nested semantic code whose distortion is low across deployment rates.
Semantic neighborhood geometry provides an operational variational surrogate for the information lost by each prefix, enabling direct optimization and measurement of successive semantic refinement.
}
\end{quote}


% =========================================================
\section{Target Figures}
% =========================================================

\paragraph{Figure 1: Rotation motivation.}
Show two full-dimensional representations with identical full-space geometry but different low-dimensional prefix quality after orthogonal rotation.

\paragraph{Figure 2: Method overview.}
Illustrate
\begin{equation}
    Z_{16}
    \subset
    Z_{32}
    \subset
    Z_{64}
    \subset
    \cdots
    \subset
    Z_D
\end{equation}
together with a full-dimensional semantic teacher distribution and prefix semantic decoders.

\paragraph{Figure 3: Semantic distortion-rate curves.}
Plot
\begin{equation}
    \cD_k^{\mathrm{ref}}
\end{equation}
against embedding storage rate for MRL, MIPIC, ESE, SDR-MRL, and independent low-dimensional models.

\paragraph{Figure 4: Price of nestedness.}
Plot
\begin{equation}
    P_{\mathrm{nest}}(d)
\end{equation}
as a function of dimension.

\paragraph{Figure 5: Refinement gain versus downstream gain.}
Scatter
\begin{equation}
    \cD_{k-1}-\cD_k
\end{equation}
against
\begin{equation}
    Q_k-Q_{k-1}
\end{equation}
across models and datasets.


% =========================================================
\section{Proposed Paper-Level Claims}
% =========================================================

If supported by the experiments, the final manuscript should make no more than the following three primary claims:

\paragraph{Claim 1: Conceptual.}
Nested dimensionality alone does not specify nested semantic information.
A Matryoshka embedding can be modeled as a multi-rate semantic representation whose quality is characterized by semantic distortion at each deployment rate.

\paragraph{Claim 2: Theoretical.}
Neighborhood-distribution distortion is a variational upper bound on conditional semantic information loss.
At the optimal semantic decoder, the distortion reduction from an added Matryoshka block equals that block's conditional semantic information.

\paragraph{Claim 3: Empirical.}
Optimizing semantic neighborhood distortion across deployment rates improves low-dimensional semantic geometry and reduces the price of nestedness relative to existing Matryoshka training objectives.

These claims are intentionally narrower than asserting a complete neural rate-distortion theory.
The proposed storage rate is operational rather than Shannon-optimal, and the semantic teacher defines a task-dependent notion of relevant information.


% =========================================================
\section{Conclusion}
% =========================================================

This proposal reframes Matryoshka Representation Learning as a semantic successive-refinement problem.
The core idea is to move from matching representations to measuring the semantic information that remains recoverable at each embedding budget.

A full-dimensional teacher induces a probabilistic semantic neighborhood.
Each prefix attempts to reconstruct this neighborhood using only its truncated geometry.
The resulting distortion is not merely a heuristic geometric alignment loss: it decomposes into conditional semantic information loss and a variational approximation gap.
This provides a direct mathematical link between local representation geometry and information preservation.

The framework suggests both a practical training objective and a richer evaluation protocol.
Rather than evaluating only task accuracy at isolated dimensions, the proposed experiments measure distortion-rate profiles, neighborhood preservation, marginal semantic refinement, and the price paid for requiring multiple capacities to coexist within a single ordered representation.

If the empirical results support the central hypothesis, SDR-MRL would provide a principled answer to a question left implicit in existing Matryoshka learning:
\begin{quote}
\emph{
what semantic information should survive first when an embedding is progressively compressed?
}
\end{quote}


% =========================================================
% Bibliography
% =========================================================

\begin{thebibliography}{99}

\bibitem{kusupati2022matryoshka}
Aditya Kusupati, Gantavya Bhatt, Aniket Rege, Matthew Wallingford,
Aditya Sinha, Vivek Ramanujan, William Howard-Snyder, Kaifeng Chen,
Sham Kakade, Prateek Jain, and Ali Farhadi.
\newblock Matryoshka Representation Learning.
\newblock In \emph{Advances in Neural Information Processing Systems},
2022.

\bibitem{huy2026mipic}
Phung Gia Huy, Hai An Vu, Minh-Phuc Truong, Thang Duc Tran,
Linh Ngo Van, Thanh Hong Nguyen, and Trung Le.
\newblock MIPIC: Matryoshka Representation Learning via Self-Distilled
Intra-Relational and Progressive Information Chaining.
\newblock In \emph{Findings of the Association for Computational
Linguistics: ACL}, 2026.

\bibitem{li2025ese}
Xianming Li, Zongxi Li, Jing Li, Haoran Xie, and Qing Li.
\newblock ESE: Espresso Sentence Embeddings.
\newblock In \emph{International Conference on Learning Representations},
2025.

\bibitem{yoon2024matryoshka}
Jinsung Yoon, Rajarishi Sinha, Sercan O. Arik, and Tomas Pfister.
\newblock Matryoshka-Adaptor: Unsupervised and Supervised Tuning for
Smaller Embedding Dimensions.
\newblock In \emph{Proceedings of the 2024 Conference on Empirical
Methods in Natural Language Processing}, 2024.

\bibitem{zhang2025smec}
Biao Zhang, Lixin Chen, Tong Liu, and Bo Zheng.
\newblock SMEC: Rethinking Matryoshka Representation Learning for
Retrieval Embedding Compression.
\newblock In \emph{Proceedings of the 2025 Conference on Empirical
Methods in Natural Language Processing}, 2025.

\bibitem{kornblith2019similarity}
Simon Kornblith, Mohammad Norouzi, Honglak Lee, and Geoffrey Hinton.
\newblock Similarity of Neural Network Representations Revisited.
\newblock In \emph{Proceedings of the 36th International Conference on
Machine Learning}, pages 3519--3529, 2019.

\bibitem{wang2020alignment}
Tongzhou Wang and Phillip Isola.
\newblock Understanding Contrastive Representation Learning through
Alignment and Uniformity on the Hypersphere.
\newblock In \emph{Proceedings of the 37th International Conference on
Machine Learning}, pages 9929--9939, 2020.

\bibitem{alemi2017vib}
Alexander A. Alemi, Ian Fischer, Joshua V. Dillon, and Kevin Murphy.
\newblock Deep Variational Information Bottleneck.
\newblock In \emph{International Conference on Learning Representations},
2017.

\bibitem{yu2020mcr2}
Yaodong Yu, Kwan Ho Ryan Chan, Chong You, Chaobing Song, and Yi Ma.
\newblock Learning Diverse and Discriminative Representations via the
Principle of Maximal Coding Rate Reduction.
\newblock In \emph{Advances in Neural Information Processing Systems},
2020.

\bibitem{equitz1991successive}
William H. R. Equitz and Thomas M. Cover.
\newblock Successive Refinement of Information.
\newblock \emph{IEEE Transactions on Information Theory},
37(2):269--275, 1991.

\end{thebibliography}

\end{document}

